\documentclass[11pt]{article}
\usepackage[margin=1in]{geometry}
\usepackage[T1]{fontenc}
\usepackage{lmodern}
\usepackage{amsmath,amssymb,amsthm}
\usepackage[hidelinks]{hyperref}
\usepackage{microtype}
\raggedright

\newtheorem*{answer}{Answer}

\begin{document}

\begin{center}
{\Large Literature-implied answer for arXiv:1703.02919}\\[4pt]
{\large Exact simultaneous bisection follows from Lyapunov convexity}
\end{center}

\paragraph{Status.}
\textbf{literature\_implied\_answer (full affirmative answer to Remark 3.4).}
This is not a new theorem produced by the run.  It is a direct specialization
of the classical finite-dimensional Lyapunov convexity theorem.

\paragraph{Source question.}
In Rainis Haller, Paavo Kuuseok, and M\"art P\~oldvere,
\emph{On convex combinations of slices of the unit ball in Banach spaces},
arXiv:1703.02919v2, Remark~3.4 on PDF page~12 follows Lemma~3.2.  Given
atomless finite nonnegative measures $\mu_1,\ldots,\mu_n$ on a
$\sigma$-algebra and $E$ measurable, the lemma finds a partition
$E=A\mathbin{\dot\cup}B$ satisfying
\[
  |\mu_i(A)-\mu_i(B)|\leq\varepsilon\qquad(1\leq i\leq n).
\]
The remark asks whether one may take $\varepsilon=0$ when $n\geq2$.

\paragraph{Supporting literature.}
Lyapunov's convexity theorem says that the range of a finite vector measure in
a Euclidean space is compact and is convex if the vector measure is atomless.
This formulation appears explicitly in the abstract and introduction on PDF
page~1 of Peng Dai and Eugene A. Feinberg,
\emph{Extension of Lyapunov's Convexity Theorem to Subranges},
arXiv:1102.2534.  Their introduction defines a vector measure
$(\mu_1,\ldots,\mu_m)$ to be atomless when every coordinate measure is
atomless, exactly matching the hypotheses of Lemma~3.2.

\paragraph{Identification.}
Restrict the measures to $E$ and define the finite-dimensional vector measure
\[
  \nu(C):=(\mu_1(C),\ldots,\mu_n(C)),
  \qquad C\subseteq E,\ C\text{ measurable}.
\]
Every coordinate is atomless, so Lyapunov's theorem makes the range
\[
  R_\nu(E):=\{\nu(C): C\subseteq E,\ C\text{ measurable}\}
\]
convex.  It contains both $0=\nu(\varnothing)$ and $\nu(E)$.  Hence it
contains their midpoint $\frac12\nu(E)$.  Choose measurable $A\subseteq E$
with $\nu(A)=\frac12\nu(E)$ and put $B=E\setminus A$.  Then, for every
$i$,
\[
  \mu_i(A)=\tfrac12\mu_i(E)=\mu_i(B).
\]

\begin{answer}
Yes.  For every finite family of atomless finite nonnegative measures and
every measurable $E$, there is an exact simultaneous bisection
$E=A\mathbin{\dot\cup}B$ with $\mu_i(A)=\mu_i(B)$ for all $i$.
\end{answer}

\paragraph{Provenance and scope.}
The supporting statement predates the source question by six years.  Bounded
search of the run indexes, exact source phrase, paper title together with
``Lyapunov'', and lemma/author combinations found no later paper explicitly
identifying Remark~3.4.  The answer is therefore recorded as
literature-implied rather than as a new proof or an explicitly announced
later answer.  This packet concerns only the measure-theoretic remark.

\begin{thebibliography}{9}
\bibitem{Haller1703}
R. Haller, P. Kuuseok, and M. P\~oldvere,
\emph{On convex combinations of slices of the unit ball in Banach spaces},
arXiv:1703.02919v2 (2017).

\bibitem{Dai1102}
P. Dai and E. A. Feinberg,
\emph{Extension of Lyapunov's Convexity Theorem to Subranges},
arXiv:1102.2534 (2011).

\bibitem{Lyapunov1940}
A. A. Lyapunov,
\emph{Sur les fonctions-vecteurs compl\`etement additives},
Izv. Akad. Nauk SSSR Ser. Mat. 4 (1940), 465--478.
\end{thebibliography}

\end{document}
